% =====================================================================
%  preambolo.tex  -  Dispense di Fisica 2
%
%  Uso in ogni dispensa (e nel testo unico):
%      \documentclass[11pt,a4paper]{article}
%      % =====================================================================
%  preambolo.tex  -  Dispense di Fisica 2
%
%  Uso in ogni dispensa (e nel testo unico):
%      \documentclass[11pt,a4paper]{article}
%      % =====================================================================
%  preambolo.tex  -  Dispense di Fisica 2
%
%  Uso in ogni dispensa (e nel testo unico):
%      \documentclass[11pt,a4paper]{article}
%      % =====================================================================
%  preambolo.tex  -  Dispense di Fisica 2
%
%  Uso in ogni dispensa (e nel testo unico):
%      \documentclass[11pt,a4paper]{article}
%      \input{preambolo}
%      \begin{document} ... \end{document}
%
%  Compatibile con pdfLaTeX, XeLaTeX e LuaLaTeX (il motore viene
%  riconosciuto da solo). Le macro e gli ambienti di tutte le
%  dispense precedenti sono qui riuniti senza conflitti.
% =====================================================================

% ---------------------------------------------------------------------
%  1. Motore, codifica e font
% ---------------------------------------------------------------------
\RequirePackage{iftex}
\ifPDFTeX
  \usepackage[utf8]{inputenc}
  \usepackage[T1]{fontenc}
\else
  \usepackage{fontspec}          % XeLaTeX / LuaLaTeX
\fi
\usepackage{lmodern}
\usepackage[italian]{babel}
\usepackage{csquotes}            % \enquote{...}, virgolette corrette

% ---------------------------------------------------------------------
%  2. Matematica
% ---------------------------------------------------------------------
\usepackage{amsmath,amssymb,mathtools}

% ---------------------------------------------------------------------
%  3. Pagina, interlinea, tipografia
% ---------------------------------------------------------------------
\usepackage{microtype}
\usepackage[a4paper,margin=2.6cm]{geometry}
\usepackage{setspace}
\usepackage{needspace}
\usepackage{graphicx}

% Interlinea: lasciare attivo un solo comando
\doublespacing
%\onehalfspacing
%\setstretch{1.25}

% Paragrafi e tolleranze di composizione
\setlength{\parindent}{1.4em}
\setlength{\parskip}{0.35ex plus 0.2ex}
\setlength{\emergencystretch}{2em}

% ---------------------------------------------------------------------
%  4. Colori, riquadri e grafica
% ---------------------------------------------------------------------
\usepackage[most]{tcolorbox}     % carica anche xcolor e tikz
\usepackage{tikz}
\usetikzlibrary{arrows.meta, calc, backgrounds}
\usetikzlibrary{decorations.markings, decorations.pathreplacing}
\usetikzlibrary{3d}

% ---------------------------------------------------------------------
%  5. Notazione fisica e matematica
% ---------------------------------------------------------------------
\newcommand{\dd}{\mathrm{d}}                          % differenziale
\newcommand{\nab}{\vec{\nabla}}                       % nabla vettoriale
\newcommand{\un}[1]{\,\mathrm{#1}}                    % unita' di misura: 5\un{m}
\DeclareMathOperator{\grad}{grad}
\DeclareMathOperator{\Div}{div}
\DeclareMathOperator{\rot}{rot}

\newcommand{\epz}{\varepsilon_0}                      % permittivita' del vuoto
\newcommand{\Ev}{\vec{E}}                             % campo elettrico
\newcommand{\Fv}{\vec{F}}                             % forza
\newcommand{\uv}{\hat{u}}                             % versore generico
\newcommand{\ur}{\ensuremath{\hat{u}_r}}              % versore radiale
\newcommand{\kc}{\ensuremath{\frac{1}{4\pi\epz}}}     % 1/(4 pi eps_0)

% ---------------------------------------------------------------------
%  6. Titoli di paragrafo (maiuscoletto, senza grassetto ne' corsivo)
% ---------------------------------------------------------------------
\newcommand{\titoletto}[1]{%
  \par\Needspace{4\baselineskip}\medskip
  \noindent{\large\scshape #1}\par\nopagebreak\smallskip\noindent}

% ---------------------------------------------------------------------
%  7. Riquadri di richiamo
%     \begin{richiamo}{titolo} ... \end{richiamo}
%     \begin{riquadro}{titolo} ... \end{riquadro}   (identico, alias)
%     Per anteporre un prefisso fisso ai titoli:
%         \renewcommand{\richiamoprefix}{Richiamo di Fisica 1 --- }
% ---------------------------------------------------------------------
\newcommand{\richiamoprefix}{}

\tcbset{richiamo style/.style={
  enhanced, breakable,
  colback=gray!5, colframe=gray!55!black,
  boxrule=0.6pt, arc=2pt,
  left=8pt, right=8pt, top=6pt, bottom=6pt,
  before skip=10pt, after skip=10pt,
  fonttitle=\normalfont\large, coltitle=white,
  before upper={\singlespacing}
}}

\newtcolorbox{richiamo}[1]{richiamo style, title={\richiamoprefix #1}}
\newtcolorbox{riquadro}[1]{richiamo style, title={\richiamoprefix #1}}

% ---------------------------------------------------------------------
%  8. Collegamenti ipertestuali (sempre per ultimo)
% ---------------------------------------------------------------------
\usepackage[hidelinks]{hyperref}

%      \begin{document} ... \end{document}
%
%  Compatibile con pdfLaTeX, XeLaTeX e LuaLaTeX (il motore viene
%  riconosciuto da solo). Le macro e gli ambienti di tutte le
%  dispense precedenti sono qui riuniti senza conflitti.
% =====================================================================

% ---------------------------------------------------------------------
%  1. Motore, codifica e font
% ---------------------------------------------------------------------
\RequirePackage{iftex}
\ifPDFTeX
  \usepackage[utf8]{inputenc}
  \usepackage[T1]{fontenc}
\else
  \usepackage{fontspec}          % XeLaTeX / LuaLaTeX
\fi
\usepackage{lmodern}
\usepackage[italian]{babel}
\usepackage{csquotes}            % \enquote{...}, virgolette corrette

% ---------------------------------------------------------------------
%  2. Matematica
% ---------------------------------------------------------------------
\usepackage{amsmath,amssymb,mathtools}

% ---------------------------------------------------------------------
%  3. Pagina, interlinea, tipografia
% ---------------------------------------------------------------------
\usepackage{microtype}
\usepackage[a4paper,margin=2.6cm]{geometry}
\usepackage{setspace}
\usepackage{needspace}
\usepackage{graphicx}

% Interlinea: lasciare attivo un solo comando
\doublespacing
%\onehalfspacing
%\setstretch{1.25}

% Paragrafi e tolleranze di composizione
\setlength{\parindent}{1.4em}
\setlength{\parskip}{0.35ex plus 0.2ex}
\setlength{\emergencystretch}{2em}

% ---------------------------------------------------------------------
%  4. Colori, riquadri e grafica
% ---------------------------------------------------------------------
\usepackage[most]{tcolorbox}     % carica anche xcolor e tikz
\usepackage{tikz}
\usetikzlibrary{arrows.meta, calc, backgrounds}
\usetikzlibrary{decorations.markings, decorations.pathreplacing}
\usetikzlibrary{3d}

% ---------------------------------------------------------------------
%  5. Notazione fisica e matematica
% ---------------------------------------------------------------------
\newcommand{\dd}{\mathrm{d}}                          % differenziale
\newcommand{\nab}{\vec{\nabla}}                       % nabla vettoriale
\newcommand{\un}[1]{\,\mathrm{#1}}                    % unita' di misura: 5\un{m}
\DeclareMathOperator{\grad}{grad}
\DeclareMathOperator{\Div}{div}
\DeclareMathOperator{\rot}{rot}

\newcommand{\epz}{\varepsilon_0}                      % permittivita' del vuoto
\newcommand{\Ev}{\vec{E}}                             % campo elettrico
\newcommand{\Fv}{\vec{F}}                             % forza
\newcommand{\uv}{\hat{u}}                             % versore generico
\newcommand{\ur}{\ensuremath{\hat{u}_r}}              % versore radiale
\newcommand{\kc}{\ensuremath{\frac{1}{4\pi\epz}}}     % 1/(4 pi eps_0)

% ---------------------------------------------------------------------
%  6. Titoli di paragrafo (maiuscoletto, senza grassetto ne' corsivo)
% ---------------------------------------------------------------------
\newcommand{\titoletto}[1]{%
  \par\Needspace{4\baselineskip}\medskip
  \noindent{\large\scshape #1}\par\nopagebreak\smallskip\noindent}

% ---------------------------------------------------------------------
%  7. Riquadri di richiamo
%     \begin{richiamo}{titolo} ... \end{richiamo}
%     \begin{riquadro}{titolo} ... \end{riquadro}   (identico, alias)
%     Per anteporre un prefisso fisso ai titoli:
%         \renewcommand{\richiamoprefix}{Richiamo di Fisica 1 --- }
% ---------------------------------------------------------------------
\newcommand{\richiamoprefix}{}

\tcbset{richiamo style/.style={
  enhanced, breakable,
  colback=gray!5, colframe=gray!55!black,
  boxrule=0.6pt, arc=2pt,
  left=8pt, right=8pt, top=6pt, bottom=6pt,
  before skip=10pt, after skip=10pt,
  fonttitle=\normalfont\large, coltitle=white,
  before upper={\singlespacing}
}}

\newtcolorbox{richiamo}[1]{richiamo style, title={\richiamoprefix #1}}
\newtcolorbox{riquadro}[1]{richiamo style, title={\richiamoprefix #1}}

% ---------------------------------------------------------------------
%  8. Collegamenti ipertestuali (sempre per ultimo)
% ---------------------------------------------------------------------
\usepackage[hidelinks]{hyperref}

%      \begin{document} ... \end{document}
%
%  Compatibile con pdfLaTeX, XeLaTeX e LuaLaTeX (il motore viene
%  riconosciuto da solo). Le macro e gli ambienti di tutte le
%  dispense precedenti sono qui riuniti senza conflitti.
% =====================================================================

% ---------------------------------------------------------------------
%  1. Motore, codifica e font
% ---------------------------------------------------------------------
\RequirePackage{iftex}
\ifPDFTeX
  \usepackage[utf8]{inputenc}
  \usepackage[T1]{fontenc}
\else
  \usepackage{fontspec}          % XeLaTeX / LuaLaTeX
\fi
\usepackage{lmodern}
\usepackage[italian]{babel}
\usepackage{csquotes}            % \enquote{...}, virgolette corrette

% ---------------------------------------------------------------------
%  2. Matematica
% ---------------------------------------------------------------------
\usepackage{amsmath,amssymb,mathtools}

% ---------------------------------------------------------------------
%  3. Pagina, interlinea, tipografia
% ---------------------------------------------------------------------
\usepackage{microtype}
\usepackage[a4paper,margin=2.6cm]{geometry}
\usepackage{setspace}
\usepackage{needspace}
\usepackage{graphicx}

% Interlinea: lasciare attivo un solo comando
\doublespacing
%\onehalfspacing
%\setstretch{1.25}

% Paragrafi e tolleranze di composizione
\setlength{\parindent}{1.4em}
\setlength{\parskip}{0.35ex plus 0.2ex}
\setlength{\emergencystretch}{2em}

% ---------------------------------------------------------------------
%  4. Colori, riquadri e grafica
% ---------------------------------------------------------------------
\usepackage[most]{tcolorbox}     % carica anche xcolor e tikz
\usepackage{tikz}
\usetikzlibrary{arrows.meta, calc, backgrounds}
\usetikzlibrary{decorations.markings, decorations.pathreplacing}
\usetikzlibrary{3d}

% ---------------------------------------------------------------------
%  5. Notazione fisica e matematica
% ---------------------------------------------------------------------
\newcommand{\dd}{\mathrm{d}}                          % differenziale
\newcommand{\nab}{\vec{\nabla}}                       % nabla vettoriale
\newcommand{\un}[1]{\,\mathrm{#1}}                    % unita' di misura: 5\un{m}
\DeclareMathOperator{\grad}{grad}
\DeclareMathOperator{\Div}{div}
\DeclareMathOperator{\rot}{rot}

\newcommand{\epz}{\varepsilon_0}                      % permittivita' del vuoto
\newcommand{\Ev}{\vec{E}}                             % campo elettrico
\newcommand{\Fv}{\vec{F}}                             % forza
\newcommand{\uv}{\hat{u}}                             % versore generico
\newcommand{\ur}{\ensuremath{\hat{u}_r}}              % versore radiale
\newcommand{\kc}{\ensuremath{\frac{1}{4\pi\epz}}}     % 1/(4 pi eps_0)

% ---------------------------------------------------------------------
%  6. Titoli di paragrafo (maiuscoletto, senza grassetto ne' corsivo)
% ---------------------------------------------------------------------
\newcommand{\titoletto}[1]{%
  \par\Needspace{4\baselineskip}\medskip
  \noindent{\large\scshape #1}\par\nopagebreak\smallskip\noindent}

% ---------------------------------------------------------------------
%  7. Riquadri di richiamo
%     \begin{richiamo}{titolo} ... \end{richiamo}
%     \begin{riquadro}{titolo} ... \end{riquadro}   (identico, alias)
%     Per anteporre un prefisso fisso ai titoli:
%         \renewcommand{\richiamoprefix}{Richiamo di Fisica 1 --- }
% ---------------------------------------------------------------------
\newcommand{\richiamoprefix}{}

\tcbset{richiamo style/.style={
  enhanced, breakable,
  colback=gray!5, colframe=gray!55!black,
  boxrule=0.6pt, arc=2pt,
  left=8pt, right=8pt, top=6pt, bottom=6pt,
  before skip=10pt, after skip=10pt,
  fonttitle=\normalfont\large, coltitle=white,
  before upper={\singlespacing}
}}

\newtcolorbox{richiamo}[1]{richiamo style, title={\richiamoprefix #1}}
\newtcolorbox{riquadro}[1]{richiamo style, title={\richiamoprefix #1}}

% ---------------------------------------------------------------------
%  8. Collegamenti ipertestuali (sempre per ultimo)
% ---------------------------------------------------------------------
\usepackage[hidelinks]{hyperref}

%      \begin{document} ... \end{document}
%
%  Compatibile con pdfLaTeX, XeLaTeX e LuaLaTeX (il motore viene
%  riconosciuto da solo). Le macro e gli ambienti di tutte le
%  dispense precedenti sono qui riuniti senza conflitti.
% =====================================================================

% ---------------------------------------------------------------------
%  1. Motore, codifica e font
% ---------------------------------------------------------------------
\RequirePackage{iftex}
\ifPDFTeX
  \usepackage[utf8]{inputenc}
  \usepackage[T1]{fontenc}
\else
  \usepackage{fontspec}          % XeLaTeX / LuaLaTeX
\fi
\usepackage{lmodern}
\usepackage[italian]{babel}
\usepackage{csquotes}            % \enquote{...}, virgolette corrette

% ---------------------------------------------------------------------
%  2. Matematica
% ---------------------------------------------------------------------
\usepackage{amsmath,amssymb,mathtools}

% ---------------------------------------------------------------------
%  3. Pagina, interlinea, tipografia
% ---------------------------------------------------------------------
\usepackage{microtype}
\usepackage[a4paper,margin=2.6cm]{geometry}
\usepackage{setspace}
\usepackage{needspace}
\usepackage{graphicx}

% Interlinea: lasciare attivo un solo comando
\doublespacing
%\onehalfspacing
%\setstretch{1.25}

% Paragrafi e tolleranze di composizione
\setlength{\parindent}{1.4em}
\setlength{\parskip}{0.35ex plus 0.2ex}
\setlength{\emergencystretch}{2em}

% ---------------------------------------------------------------------
%  4. Colori, riquadri e grafica
% ---------------------------------------------------------------------
\usepackage[most]{tcolorbox}     % carica anche xcolor e tikz
\usepackage{tikz}
\usetikzlibrary{arrows.meta, calc, backgrounds}
\usetikzlibrary{decorations.markings, decorations.pathreplacing}
\usetikzlibrary{3d}

% ---------------------------------------------------------------------
%  5. Notazione fisica e matematica
% ---------------------------------------------------------------------
\newcommand{\dd}{\mathrm{d}}                          % differenziale
\newcommand{\nab}{\vec{\nabla}}                       % nabla vettoriale
\newcommand{\un}[1]{\,\mathrm{#1}}                    % unita' di misura: 5\un{m}
\DeclareMathOperator{\grad}{grad}
\DeclareMathOperator{\Div}{div}
\DeclareMathOperator{\rot}{rot}

\newcommand{\epz}{\varepsilon_0}                      % permittivita' del vuoto
\newcommand{\Ev}{\vec{E}}                             % campo elettrico
\newcommand{\Fv}{\vec{F}}                             % forza
\newcommand{\uv}{\hat{u}}                             % versore generico
\newcommand{\ur}{\ensuremath{\hat{u}_r}}              % versore radiale
\newcommand{\kc}{\ensuremath{\frac{1}{4\pi\epz}}}     % 1/(4 pi eps_0)

% ---------------------------------------------------------------------
%  6. Titoli di paragrafo (maiuscoletto, senza grassetto ne' corsivo)
% ---------------------------------------------------------------------
\newcommand{\titoletto}[1]{%
  \par\Needspace{4\baselineskip}\medskip
  \noindent{\large\scshape #1}\par\nopagebreak\smallskip\noindent}

% ---------------------------------------------------------------------
%  7. Riquadri di richiamo
%     \begin{richiamo}{titolo} ... \end{richiamo}
%     \begin{riquadro}{titolo} ... \end{riquadro}   (identico, alias)
%     Per anteporre un prefisso fisso ai titoli:
%         \renewcommand{\richiamoprefix}{Richiamo di Fisica 1 --- }
% ---------------------------------------------------------------------
\newcommand{\richiamoprefix}{}

\tcbset{richiamo style/.style={
  enhanced, breakable,
  colback=gray!5, colframe=gray!55!black,
  boxrule=0.6pt, arc=2pt,
  left=8pt, right=8pt, top=6pt, bottom=6pt,
  before skip=10pt, after skip=10pt,
  fonttitle=\normalfont\large, coltitle=white,
  before upper={\singlespacing}
}}

\newtcolorbox{richiamo}[1]{richiamo style, title={\richiamoprefix #1}}
\newtcolorbox{riquadro}[1]{richiamo style, title={\richiamoprefix #1}}

% ---------------------------------------------------------------------
%  8. Collegamenti ipertestuali (sempre per ultimo)
% ---------------------------------------------------------------------
\usepackage[hidelinks]{hyperref}
